\section{Conclusion}
\label{sec:conclusion}

We have presented an integrated three-stage deep-learning framework for pressurized water reactor fault diagnosis, combining semi-supervised anomaly detection, interpretable 18-class accident classification, and physics-informed trajectory prediction into a single pipeline evaluated on the \nppad{} dataset. Each stage performs competitively on its own---F1 $\approx 0.9995$ for anomaly screening, $71.2\%$ accuracy for 18-class accident classification, and $68.6\%$/$39.6\%$ reductions in physics and conservation residuals for the physics-informed surrogate relative to a data-only baseline---but the central contribution of the paper is the composition. The physics-informed surrogate provides a physics-grounded consistency check on the data-driven classifier, and the SHAP attribution layer provides a mechanistically auditable explanation for each classification, together yielding a diagnostic pipeline whose outputs can be inspected at every stage by both software and human reviewers.

The framework is deliberately modular: the physics terms of the surrogate, the class taxonomy of the classifier, and the sensor channels of the screening stage can each be adapted to other reactor designs without changing the overall pipeline structure. This makes the approach a practical candidate for the monitoring--diagnosis--prognosis backbone of digital-twin systems for current-fleet PWRs and for upcoming advanced reactor designs alike.

Future work falls along three tracks. First, we will replace the deterministic classifier and surrogate with uncertainty-aware variants (deep ensembles or approximate Bayesian methods) so that the physics-consistency score can be interpreted jointly with predictive uncertainty. Second, we will extend the evaluation to severity-stratified training/test splits to directly measure the extrapolation benefit conferred by the physics-informed loss. Third, we will explore domain-adaptation techniques that bridge the simulator-to-plant gap, which is the principal remaining obstacle to deployment on real instrument streams.
